% Options for packages loaded elsewhere
\PassOptionsToPackage{unicode}{hyperref}
\PassOptionsToPackage{hyphens}{url}
\PassOptionsToPackage{dvipsnames,svgnames,x11names}{xcolor}
\documentclass[
  10pt,
]{ctexart}
\usepackage{xcolor}
\usepackage[margin=2cm]{geometry}
\usepackage{amsmath,amssymb}
\setcounter{secnumdepth}{-\maxdimen} % remove section numbering
\usepackage{iftex}
\ifPDFTeX
  \usepackage[T1]{fontenc}
  \usepackage[utf8]{inputenc}
  \usepackage{textcomp} % provide euro and other symbols
\else % if luatex or xetex
  \usepackage{unicode-math} % this also loads fontspec
  \defaultfontfeatures{Scale=MatchLowercase}
  \defaultfontfeatures[\rmfamily]{Ligatures=TeX,Scale=1}
\fi
\usepackage{lmodern}
\ifPDFTeX\else
  % xetex/luatex font selection
  \setmainfont[]{DejaVu Sans}
\fi
% Use upquote if available, for straight quotes in verbatim environments
\IfFileExists{upquote.sty}{\usepackage{upquote}}{}
\IfFileExists{microtype.sty}{% use microtype if available
  \usepackage[]{microtype}
  \UseMicrotypeSet[protrusion]{basicmath} % disable protrusion for tt fonts
}{}
\makeatletter
\@ifundefined{KOMAClassName}{% if non-KOMA class
  \IfFileExists{parskip.sty}{%
    \usepackage{parskip}
  }{% else
    \setlength{\parindent}{0pt}
    \setlength{\parskip}{6pt plus 2pt minus 1pt}}
}{% if KOMA class
  \KOMAoptions{parskip=half}}
\makeatother
\setlength{\emergencystretch}{3em} % prevent overfull lines
\providecommand{\tightlist}{%
  \setlength{\itemsep}{0pt}\setlength{\parskip}{0pt}}
\usepackage{bookmark}
\IfFileExists{xurl.sty}{\usepackage{xurl}}{} % add URL line breaks if available
\urlstyle{same}
% fallback for those not using the hyperref driver hyperxmp:
\makeatletter
\@ifundefined{xmpquote}{\newcommand{\xmpquote}[1]{#1}}{}
\makeatother
\hypersetup{
  colorlinks=true,
  linkcolor={Maroon},
  filecolor={Maroon},
  citecolor={Blue},
  urlcolor={Blue},
  pdfcreator={LaTeX via pandoc}}

\author{}
\date{}

\usepackage{pdflscape,fvextra}
\setlength{\tabcolsep}{3pt}
\begin{document}

\section{Prospective method-selection stress test
(v2)}\label{prospective-method-selection-stress-test-v2}

Status: input and analysis plan frozen at commit \texttt{8b14325}; 15
prescore DeepSeek choices were frozen at \texttt{bbf5fb7} and
provider-visible trace grading passed 15/15. Choice JSON SHA-256:
\texttt{a87cf337fb8a74\allowbreak{}e943dc6237a6a0\allowbreak{}ee629d53aa3c0e\allowbreak{}ee1e4b0f41cf1e\allowbreak{}5b757949}.
The 15 choices select B2 five times, B1k five times, B1 three times and
B0p twice. No partition benchmark outcome was computed before this
freeze. The original five-partition run subsequently \textbf{failed} at
partition 02 because its sole \texttt{-1} rating cannot be stratified by
the official \texttt{random\_\allowbreak{}simple} splitter. This is a
within-TRANSCRIPT stress test, \textbf{not} an independent
external-dataset validation or evidence of improved drug efficacy.

Five disjoint sets of disease IDs were deterministically derived from
checksum-pinned TRANSCRIPT v2.0.0. The committed
\texttt{configs/\allowbreak{}method\_\allowbreak{}selection\_\allowbreak{}partition\_\allowbreak{}v2.\allowbreak{}json} fixes the data
manifest, four eligible methods (B0p, B1k, B1, B2),
\texttt{random\_\allowbreak{}simple} protocol, 20 outer seeds, five-fold model
selection and three independent model choices per partition.
\texttt{configs/\allowbreak{}method\_\allowbreak{}selection\_\allowbreak{}partition\_\allowbreak{}v2\_\allowbreak{}cases.\allowbreak{}json} contains
only label-blind input summaries and staged-data paths; it was committed
before model selection. Disease groups are disjoint, but drugs and the
parent dataset are shared, and outer test items within a group are
reused across runs. Do not call the resulting 15 choices 15 independent
biomedical validation cohorts.

Run order is binding:

\begin{enumerate}
\def\labelenumi{\arabic{enumi}.}
\tightlist
\item
  Verify the source checksum and stage partition data with
  \texttt{python\ -m\ scripts.\allowbreak{}stage\_\allowbreak{}method\_\allowbreak{}selection\_\allowbreak{}partition\_\allowbreak{}v2}.
  Do not inspect partition outcomes before choices are frozen.
\item
  Set \texttt{DEEPSEEK\_\allowbreak{}API\_\allowbreak{}KEY} in the local process environment or an
  ignored repository \texttt{.\allowbreak{}env}, then run
  \texttt{python\ -m\ scripts.\allowbreak{}run\_\allowbreak{}method\_\allowbreak{}selection\_\allowbreak{}partition\_\allowbreak{}v2\ -\/\allowbreak{}-recover-attempt-trace\ artifacts/\allowbreak{}reports/\allowbreak{}traces/\allowbreak{}20260926T152907345858Z\_\allowbreak{}method\_\allowbreak{}selection\_\allowbreak{}v2\_\allowbreak{}99b0cdaf25da4b\allowbreak{}2d9e950ede1436\allowbreak{}c91d.\allowbreak{}jsonl}.
  This recovery revalidates all six previous provider responses
  (including the earlier recovered first choice), in order, and does not
  resample them. The runner refuses to overwrite its choice file and
  refuses to run after partition outcome directories exist. Each choice
  stores a provider-trace SHA-256. Immediately run
  \texttt{python\ -m\ evals.\allowbreak{}grade\_\allowbreak{}method\_\allowbreak{}selection\_\allowbreak{}traces} and
  require 15/15 provider-visible request/response matches before
  computing outcomes. Preserve the resulting choice JSON, its SHA-256
  and its JSONL/provider traces before running models. Never commit
  credentials or unreviewed provider traces.
\item
  Prepare the pinned RECeSS upstream clone (commit
  \texttt{a7f11077271ced\allowbreak{}f3a98a82e3dc74\allowbreak{}b6fc0e93986e}) with
  \texttt{python\ -m\ scripts.\allowbreak{}prepare\_\allowbreak{}recess\_\allowbreak{}official\_\allowbreak{}upstream}. This
  traced, idempotent step verifies the commit and applies both local
  patches with \texttt{git\ apply\ -\/\allowbreak{}-recount} (their hunk counts
  require recount). Then use
  \texttt{python\ -m\ scripts.\allowbreak{}run\_\allowbreak{}method\_\allowbreak{}selection\_\allowbreak{}partition\_\allowbreak{}v2\_\allowbreak{}benchmarks}
  for a read-only preflight. \textbf{Only after step 2 has passed 15/15
  trace grading}, use the same command with \texttt{-\/\allowbreak{}-execute} to run
  all five partitions, four methods and 20 seeds through the official
  wrappers and automatically score them. This entrypoint refuses
  existing outcome directories and will not run without the frozen
  choices. The underlying per-case wrapper commands remain documented in
  the script for manual audit.
\item
  The batch command automatically scores after all official outputs
  pass. For an independent recheck, run
  \texttt{python\ -m\ evals.\allowbreak{}score\_\allowbreak{}method\_\allowbreak{}selection\_\allowbreak{}partition\_\allowbreak{}v2\ -\/\allowbreak{}-output\ artifacts/\allowbreak{}reports/\allowbreak{}method\_\allowbreak{}selection\_\allowbreak{}partition\_\allowbreak{}v2\_\allowbreak{}regrade.\allowbreak{}json}.
  The scorer refuses any failed choice trace and verifies frozen file
  hashes, complete and identical seed order for all four methods, finite
  official NS-AUC values and exactly three choices per case. It reports
  chosen method vs fixed B2, seed-paired differences, and an oracle
  upper bound; no post-hoc best-of-three selection or treating repeats
  as independent cohorts is permitted.
\end{enumerate}

Every stage writes JSONL trace receipts under ignored
\texttt{artifacts/\allowbreak{}}. A trace records visible calls and outputs, not a
model\textquotesingle s private reasoning. The historical split-level
selector had only two decisions; the new paired audit finds +0.039825
NS-AUC over fixed B2 on the reused weak holdout and a tie on random
splits, but does not establish generalizable LLM method-selection skill.
Even a positive v2 result would require a genuinely independent dataset
and prospective replication before such a claim.

\subsubsection{Pre-outcome format-only
amendment}\label{pre-outcome-format-only-amendment}

The first live v2 call (partition 01, repeat 1) returned a valid
\texttt{B2} tool choice, but its 718-character reason exceeded the
then-local 200-character cap. The command stopped before saving a choice
report; no partition benchmark was run or viewed. The complete, redacted
and SHA-256-chained
\href{../benchmark/results/traces/method_selection_v2_initial_provider_attempt.jsonl}{provider
response} (file SHA-256
\texttt{4ed955d92ac302\allowbreak{}61bbe79233abc6\allowbreak{}bf6e6117beb472\allowbreak{}dc5950bbbd1f73\allowbreak{}4121e96e})
and
\href{../benchmark/results/traces/method_selection_v2_initial_validation_failure.jsonl}{validation-failure
trace} (SHA-256
\texttt{e0b3b62fd6e5c3\allowbreak{}a0a0e36b402a04\allowbreak{}3699b93b3a10d6\allowbreak{}d4213d2c0b1c2b\allowbreak{}325a0b3f})
were preserved before the cap changed to 1000 characters. The schema,
blind input, eligible methods, split, scoring metric and outcome plan
did not change. Recovery checks byte-equivalent model request
construction and uses the original provider-visible response as the
first of 15 choices, avoiding selective redrawing.

The next run recovered that response, then validated four more choices.
Its sixth provider response (partition 02, repeat 3) had a
1009-character reason and exceeded the amended 1000-character cap. Again
no choice report was saved and no benchmark outcome was run or viewed.
The failed attempt trace and all six provider responses were preserved
with SHA-256 chains. The reason cap was raised to 4096 characters, well
above the two observed responses while the 400-token request limit
stayed unchanged. The recovery function replays each exact prior
request/response against its frozen blind case, checks already-validated
choices and provider hashes, and continues only from choice seven. No
methods, blind input, labels, split, scoring plan or provider response
were changed or redrawn.

\subsubsection{Runner feasibility failure and disclosed
amendment}\label{runner-feasibility-failure-and-disclosed-amendment}

After the choices were frozen, the original official run completed
partition 01 then stopped at partition 02 B2. A traced reproduction
raised
\texttt{ValueError:\ The\ least\ populated\ classes\ in\ y\ have\ only\ 1\ member\ .\allowbreak{}.\allowbreak{}.\allowbreak{}\ {[}-1{]}}
inside \texttt{stanscofi.\allowbreak{}training\_\allowbreak{}testing.\allowbreak{}random\_\allowbreak{}simple\_\allowbreak{}split} before
that partition produced a score. The original
\href{../artifacts/traces/20260926T153328294433Z_method_selection_partition_v2_benchmarks_0b04ca0f58ec49a7abdc38accf0a9e8a.jsonl}{batch
failure trace} has SHA-256
\texttt{2dc668f67e4415\allowbreak{}5de29b8646510c\allowbreak{}7aea757a0745fe\allowbreak{}dcf740b6795281\allowbreak{}27824661};
the
\href{../artifacts/traces/20260926T153805204265Z_recess_official_b2_1d011824597f421992787fb203859607.jsonl}{official
wrapper failure trace} has SHA-256
\texttt{232468aa0f27c1\allowbreak{}3b918620b8f20d\allowbreak{}9a02e39cb92d03\allowbreak{}bb7573936f61cf\allowbreak{}8ac199af}.
Explicit negative counts across the five frozen groups are 0, 1, 0, 7
and 3. No NS-AUC values from partition 01 were examined to choose the
amendment.

The separately labeled \textbf{positive-only feasibility amendment} maps
every explicit \texttt{-1} rating to \texttt{0} (unknown) in all five
partitions while preserving positives and both feature tables. This is a
post-failure data-protocol change, so its eventual score must
\textbf{not} be presented as the untouched preregistered v2 result. It
retains the previously frozen 15 method choices, uses separate staged
data and output directories, and still requires identical official
runner splits, seeds and five-fold tuning within each partition across
B0p/B1k/B1/B2. The staged
\href{../artifacts/reports/method_selection_partition_v2_positive_only_manifest.json}{amendment
manifest} has SHA-256
\texttt{5b512abd97df4d\allowbreak{}e2a85ea821a10a\allowbreak{}c0e54293d06be1\allowbreak{}3c9e67cfa88abd\allowbreak{}56f7ce35}
and its read-only preflight found no amended outcome directories. It was
frozen before any amended run. In a fresh clone, first stage the
original partition data, then use
\texttt{python\ -m\ scripts.\allowbreak{}stage\_\allowbreak{}method\_\allowbreak{}selection\_\allowbreak{}partition\_\allowbreak{}v2\_\allowbreak{}positive\_\allowbreak{}only\ -\/\allowbreak{}-verify-existing-manifest}
to rebuild the amended data and verify the committed manifest without
overwriting it. Run
\texttt{python\ -m\ scripts.\allowbreak{}run\_\allowbreak{}method\_\allowbreak{}selection\_\allowbreak{}partition\_\allowbreak{}v2\_\allowbreak{}benchmarks\ -\/\allowbreak{}-positive-only-manifest\ artifacts/\allowbreak{}reports/\allowbreak{}method\_\allowbreak{}selection\_\allowbreak{}partition\_\allowbreak{}v2\_\allowbreak{}positive\_\allowbreak{}only\_\allowbreak{}manifest.\allowbreak{}json}
for read-only preflight. The scorer writes a distinct amended status and
manifest hash; \texttt{-\/\allowbreak{}-execute} is deliberately non-overwriting and
requires a fresh outcome directory.

\subsubsection{Amended result (descriptive negative
finding)}\label{amended-result-descriptive-negative-finding}

The amended official runner completed all 5 partitions × 4 methods × 20
outer seeds with 5-fold inner tuning and exact seed-order matching. The
\href{../artifacts/reports/method_selection_partition_v2_positive_only_score.json}{scored
report} has SHA-256
\texttt{76a29338e361fc\allowbreak{}5c9f6cfb96d918\allowbreak{}6c059563d1cda8\allowbreak{}a66cf28246f1aa\allowbreak{}1fcf27a3};
an independent regrade matched every outcome, source hash and score. The
frozen 15 choices passed provider-visible trace grading 15/15, but
choosing among methods from label-blind metadata \textbf{underperformed}
simply fixing B2. No selected non-B2 method beat B2 in its partition:
five B2 ties and ten losses across the 15 individual choices. The
repeats share outcomes and are not independent datasets.

\textbf{Partition 01 · B2, B1k, B1}

\begin{itemize}
\tightlist
\item
  \textbf{Selected mean NS-AUC：} 0.37087
\item
  \textbf{Fixed B2 NS-AUC：} 0.46433
\item
  \textbf{Paired delta vs B2：} -0.09346
\item
  \textbf{Seed wins/ties/losses：} 3/0/17
\end{itemize}

\textbf{Partition 02 · B2, B2, B0p}

\begin{itemize}
\tightlist
\item
  \textbf{Selected mean NS-AUC：} 0.50407
\item
  \textbf{Fixed B2 NS-AUC：} 0.50611
\item
  \textbf{Paired delta vs B2：} -0.00204
\item
  \textbf{Seed wins/ties/losses：} 12/0/8
\end{itemize}

\textbf{Partition 03 · B1, B1k, B0p}

\begin{itemize}
\tightlist
\item
  \textbf{Selected mean NS-AUC：} 0.41324
\item
  \textbf{Fixed B2 NS-AUC：} 0.52828
\item
  \textbf{Paired delta vs B2：} -0.11504
\item
  \textbf{Seed wins/ties/losses：} 2/0/18
\end{itemize}

\textbf{Partition 04 · B1k, B2, B1}

\begin{itemize}
\tightlist
\item
  \textbf{Selected mean NS-AUC：} 0.43971
\item
  \textbf{Fixed B2 NS-AUC：} 0.56251
\item
  \textbf{Paired delta vs B2：} -0.12280
\item
  \textbf{Seed wins/ties/losses：} 0/0/20
\end{itemize}

\textbf{Partition 05 · B2, B1k, B1k}

\begin{itemize}
\tightlist
\item
  \textbf{Selected mean NS-AUC：} 0.30476
\item
  \textbf{Fixed B2 NS-AUC：} 0.52298
\item
  \textbf{Paired delta vs B2：} -0.21822
\item
  \textbf{Seed wins/ties/losses：} 0/0/20
\end{itemize}

\textbf{Partition Equal-weight partition mean · ---}

\begin{itemize}
\tightlist
\item
  \textbf{Selected mean NS-AUC：} \textbf{0.40653}
\item
  \textbf{Fixed B2 NS-AUC：} \textbf{0.51684}
\item
  \textbf{Paired delta vs B2：} \textbf{-0.11031}
\item
  \textbf{Seed wins/ties/losses：} 0 positive / 5 negative partitions
\end{itemize}

This is evidence \textbf{against} this particular feature-only LLM
selector in these amended within-dataset partitions, not evidence that
LLMs can never assist with model selection. The altered rating protocol,
shared parent dataset, only five disease groups and reused seeds
preclude a clean confirmatory or clinical claim. A better next design
would let a selector use training-only inner-CV evidence under a newly
frozen protocol, then evaluate it on genuinely unseen outer partitions
or an external dataset; these five outcomes must not become that future
test set.

\end{document}
